\section{Introducere}

În societatea contemporană, influența dezvoltărilor mijloacelor de calcul este incontestabilă și se resfrânge asupra tuturor aspectelor vieții.
Acest impact semnificativ nu este unul întâmplător, ci rezultatul a decenii de inovații.
Două arii conexe care s-au remarcat timpuriu a fi semnificative în această evoluție au fost criptografia și criptanaliza. \\

Asemănător cu orice alt domeniu, pe parcursul timpului, în criptografie, au apărut o serie de provocări și puncte vulnerabile, dintre acestea evidențiindu-se: autentificarea și integritatea datelor.
Autentificarea asigură veridicitatea identității iar integritatea datelor garantează absența modificărilor nedorite.
În acest context, se remarcă funcțiile hash, un instrument criptografic cu aplicații variate.
Însăși definiția funcției hash demonstrează caracterul versatil al utilității sale: transformă un șir de lungime arbitrară într-un șir de lungime fixă.
Unidirecționalitatea funcției hash determină practicalitatea acesteia.
Începând sfărșitul celui de-al doilea mileniu, numeroase structuri de funcții hash au fost propuse și analizate.
Prima generație de funcții hash moderne a fost reprezentată de MD5~\cite{rivest1992md5}, bazată pe construcția Merkle-Damgård~\cite{merkle1989,damgard1989}.
Vulnerabilitățile acestei funcții la atacurile cu preimagine sau extinderea lungimii mesajului au devenit rapid evidente.
În consecință, au urmat două alte iterații menite să adreseze problemele MD5: SHA-1~\cite{nist1995sha1} și SHA-2~\cite{nist2015sha2}.
SHA-1 a fost ulterior compromisă, dar SHA-2 rămâne în continuare standardul la nivel de industrie.
În ciuda adoptării generalizate, SHA-2 este fundamentată pe aceeași construcție Merkele-Damgård.
Defectele structurale ale acestei construcții au fost relevate de timp.
Construcția s-a dovedit vulnerabilă la o multitidudine de atacuri, inclusiv cele cu preimagine și extinderea lungimii mesajului.
Având în vedere acest context, apariția SHA-3 a fost inevitabilă.
Structura SHA-3 se îndepărtează de structura definită de Merkle și Damgård, și adoptă construcția sponge~\cite{bertoni2007sponge}. \\

În procesul de confuzie și difuzie caracteristic funcțiilor hash~\cite{shannon1949}, se remarcă numeroase de structuri și construcții.
Majoritatea funcțiilor hash moderne se bazează pe substituții, transpoziții, permutări și operații aritmetice.
Aceste metode sunt eficiente, însă proprietățile lor criptografice nu sunt garantate de un model matematic unitar.
Mularea pe o serie de criterii statistice, fără a exista o garanție structurală produce un efect de avalanșă incert, întrucât implică o serie de ipoteze și presupuneri care, în eventualitatea unei vulnerabilități, pot fi invalidate sau propagate la dependenți.
Un instrument care conferă prin definiție confuzie și difuzie este reprezentat de sistemele dinamice haotice~\cite{alvarez2006}, o subcategorie a sistemelor dinamice discrete.
Un astfel de sistem descrie o transformare deterministă care pare complet predictibilă la o primă inspecție, însă, în realitate, manifestă un comportament complex și aparent aleator, caracterizat de o sensibilitate extremă la condițiile inițiale.
Aceste sisteme prezintă trei proprietăți definitorii care se aliniază remarcabil cu cerințele criptografice formulate de Shannon~\cite{shannon1949,devaney1989,alvarez2006}.
Prima este tranzitivitatea topolgică, iar a doua este densitatea orbitelor periodice.
Cea de a treia proprietate, sensibilitatea la condițiile inițiale, este asigurată de primele două proprietăți~\cite{banks1992}.
Astfel, faptul că securitatea sistemelor haotice nu se întemeiază pe nicio construcție euristică este esențial.
Absența unei structuri algebrice exploatabile, asociată de regulă cu vulnerabilitatea la atacurile cuantice, transformă această aparentă lipsă de structură într-un atu de securitate.

Având în vedere aceste considerente, lucrarea de față propune o nouă funcție hash bazată pe sisteme dinamice discrete haotice.
Securitatea acestei funcții hash se bazează pe proprietățile sistemelor haotice îmbinate cu tehnicile clasice de confuzie și difuzie.
Ce diferențiază această funcție hash de cele existente este structura optimizată pentru procesare paralelă.
Folosindu-se de arbori Merkle și de construcția sponge, această funcție hash are o performanță superioară în comparație cu alte funcții similare.
Valorificând sistemele hardware moderne multicore, această funcție are o viteză de procesare care o face utilizabilă în scenarii de producție. \\

În continuare, în Capitolul 2 vor fi prezentate sistemele haotice utilizate în această lucrare.
Capitolul 3 va include o descriere exhaustivă a arhitecturii funcției hash propuse, împreună cu o implementare în pseudocod a funcției.
În capitolul 4 va fi prezentată atât analiza criptografica a funcției hash, cât și analiza performanței acesteia, mai ales în comparație cu alte funcții hash similare.
În final, în capitolul 5 vor fi prezentate concluziile și viitoare direcții de dezvoltare.